\documentclass[a4paper, 11pt, fleqn, twoside]{scrreprt}
% escapeinside
%\usepackage[gray]{xcolor}
\usepackage[T1]{fontenc}        % T1-Fonts
\usepackage[utf8]{inputenc}     % Enkodierung
\usepackage[english]{babel}
\usepackage{amsmath}
\usepackage{amssymb}
%\usepackage{proof} % 
%%http://www.logicmatters.net/resources/ndexamples/proofsty.html
%\usepackage{bussproofs}
\usepackage{hyperref}
\usepackage{subcaption}
\usepackage{tikz}
\usepackage{natbib}
\usepackage{abstract}
\usepackage[automark, headsepline]{scrlayer-scrpage}

% Kapitelüberschrift in der Kopfzeile
%\usepackage[automark]{scrpage2} % Schickerer Satzspiegel mit KOMA-Script
%\pagestyle{scrheadings}

% Minted
\usepackage{minted}
\usemintedstyle[Haskell]{trac, fontsize=\small}%,frame=lines, framesep=10pt}
\setmintedinline{style = bw}

% Für schönere Tabellen (optional)
% \usepackage{booktabs}           % Netteres Tabellenlayout
% \usepackage{multicol}           % Mehrspaltige Bereiche
% \usepackage{tabularx}           % Tabellen mit fester Breite

% Für Listings
% \usepackage{listings}

% Eine kleine Hilfe für offene Lücken
\newcommand{\todo}[1]{\marginpar{\textbf{TODO:} #1}}
\newcommand{\hinl}[1]{\mintinline{Haskell}{#1}}
\newcommand{\cinl}[1]{\mintinline{Coq}{#1}}

\begin{document}

\pagenumbering{roman} % römische Seitenzahlen

\begin{titlepage}
	\vspace*{3cm}
	\centering
	{\huge\bfseries Modelling Call-Time Choice as Effect using Scoped Free 
	Monads\par}
	\vspace{1cm}
	\textbf{Niels Bunkenburg} \par 
	\vspace{6cm}
	\textbf{Master's Thesis} \par
	Programming Languages and Compiler Construction \par
	Department of Computer Science \par
	Kiel University
	\vfill
	Advised by\par
	Priv.-Doz. Dr. Frank Huch \par
	M. Sc. Sandra Dylus
	\vfill
	% Randloses Drucken nicht möglich...
	\tikz[remember picture,overlay] \node[opacity=0.3,inner sep=0pt] at 
	(9.5,-1.5){\includegraphics{img/cau-siegel-1400.png}};
	{\large \today\par}
\end{titlepage}

\newpage
\thispagestyle{empty}
\strut
\newpage

\chapter*{Erklärung der Urheberschaft}
\vspace{2cm}
Ich erkläre hiermit an Eides statt, dass ich die vorliegende Arbeit
ohne Hilfe Dritter und ohne Benutzung anderer als der angegebenen
Hilfsmittel angefertigt habe. Aus fremden Quellen direkt oder
indirekt übernommene Gedanken sind als solche kenntlich gemacht. Die
Arbeit wurde bisher in gleicher oder ähnlicher Form in keiner anderen
Prüfungsbehörde vorgelegt und auch noch nicht veröffentlicht.

\vspace{4cm}
\hspace{1cm} $\overline{~~~~~~~~~~\mbox{Ort, Datum}~~~~~~~~~~}$ \hfill $\overline{~~~~~~~~~~~~~\mbox{Unterschrift}~~~~~~~~~~~~~}$ \hspace{1cm}

\newpage
\mbox{}
\thispagestyle{empty}
\newpage

\begin{abstract}

\end{abstract}

\newpage
\strut
\thispagestyle{empty}
\newpage

% Verzeichnisse
\renewcommand{\contentsname}{Contents}
\tableofcontents   % Inhaltsverzeichnis
%\listoffigures     % Abbildungsverzeichnis
% \listoftables      % Tabellenverzeichnis
% \lstlistoflistings % Abbildungsverzeichnis

\newpage               % Expliziter Umbruch für Seitenzahlen
\pagenumbering{arabic} % arabische Seitenzahlen

% Inhalt

\chapter{Introduction}

\chapter{Preliminaries}

\section{Coq}
\label{sec:coqIntro}
\begin{itemize}
\item Introduce the necessary Coq concepts to understand the paper
\end{itemize}

\section{Haskell}
\begin{itemize}
\item Introduce the necessary Haskell concepts to understand the paper
\end{itemize}
\subsection{Monad and MonadPlus}

\section{Curry}
\begin{itemize}
\item Introduce the necessary Curry concepts to understand the paper
\end{itemize}
\subsection{Non-strictness}
\subsection{Sharing}
\subsection{Non-determinism}

\section{Modelling Curry Programs using Monadic Code Transformation}
\begin{itemize}
\item Why is the naive MonadPlus approach not sufficient to model Curry 
semantic?
\item Motivate usage of monadic data types
\item Introduce explicit sharing
\end{itemize}

Modelling Curry programs in a language like Haskell requires a transformation 
of non-deterministic code into a semantically equivalent, deterministic 
program. First, we have a look at the direct representation of non-determinism 
used in the KiCS2 implementation as described by \citet{brassel2011kics2}.

Non-determinism in Curry is not limited to \textit{flat} non-determinism but 
can occur within components of data structures and anywhere in a computation. 
This means that expressing non-determinism via Haskell's list monad is not 
sufficient to model Curry's non-determinism. \todo{Example}
Instead, existing data types receive additional constructors that represent 
failure and the choice between two values. For example, the extended list data 
type looks as follows.

\begin{minted}{Haskell}
data List a = Nil | Cons a (List a) | Choice (List a) (List a) | Fail
\end{minted}

Since this transformation adds new constructors, all functions need to cover 
these cases, too. The new rules return \hinl{Fail} if the 
function's argument is a failed computation and distribute function calls to 
both branches if the argument is a choice.

One issue with this approach is that call-time choice is not implemented yet. 
If a choice is duplicated during evaluation, this information cannot be 
recovered later. Therefore, each \hinl{Choice} constructor has 
an additional \hinl{ID} argument that identifies the same 
choices. Since each choice needs a fresh ID, functions use an additional 
\hinl{IDSupply} argument when choices are created.

The evaluation of a non-deterministic value is implemented by transforming the 
value into a search tree which can be traversed with different search 
strategies. In the process, each choice ID's decision is stored and then 
repeated if the same ID is encountered again.

While this approach is useful when the host language supports laziness and 
sharing, another approach is necessary to model these effects when they are not 
built into the language. 

\citet{fischer2009purely} introduce a monadic representation of non-determinism 
that supports sharing and non-strict evaluation. Out of simplicity, the 
implementation idea is presented in Haskell, similar to the approach of the 
original authors, using the example of permutation sort.
The algorithm consists  of three components. Firstly, a function 
\hinl{insert} that inserts an element non-deterministically at 
every possible position within a list.

\begin{minted}{Haskell}
insert :: MonadPlus m => a -> [a] -> m [a]
insert x xs = return (x:xs)
      `mplus` case xs of
                []     -> mzero
                (y:ys) -> do zs <- insert x ys
                             return (y:zs)
\end{minted}

The second part is the function \hinl{perm} that inserts the 
head of a given list into the permutations of the list's tail.
\begin{minted}{Haskell}
perm :: MonadPlus m => [a] -> m [a]
perm [] = return []
perm (x:xs) = do ys <- perm xs
                 zs <- insert x ys
                 return zs
\end{minted}

Finally, the function \hinl{sort} generates permutations and 
then tests whether they are sorted.

\begin{minted}{Haskell}
sort :: MonadPlus m => [Int] -> m [Int]
sort xs = do ys <- perm xs
             guard (isSorted ys)
             return ys
\end{minted}

The function \hinl{isSorted} compares each element in a list to 
the next one to determine whether the list is sorted. When we test this 
implementation, we can see that the runtime increases significantly when adding 
even a few elements.

\begin{minted}{Haskell}
*Test> sort [9, 8..1] :: [[Int]]
[[1,2,3,4,5,6,7,8,9]]
(0.69 secs, 717,914,784 bytes)
*Test> sort [10, 9..1] :: [[Int]]
[[1,2,3,4,5,6,7,8,9,10]]
(6.67 secs, 7,437,960,280 bytes)
*Test> sort [11, 10..1] :: [[Int]]
[[1,2,3,4,5,6,7,8,9,10,11]]
(77.54 secs, 84,743,416,080 bytes)
\end{minted}

The reason for the factorial runtime is that the implementations is needlessly 
strict. A list of length $n$  has $n!$ permutations, all of which are generated 
when running \hinl{sort}. This matches our observation above, 
since adding a tenth element increases the runtime by a factor of 10 and an 
eleventh element multiplies the runtime of the ten-element list by eleven. 

If we consider the implementation of \hinl{isSorted}, we can see 
that, as soon as the comparison of two elements yields 
\hinl{False}, the function returns \hinl{False} 
and does not evaluate the remainder of the list.

\begin{minted}{Haskell}
isSorted :: [Int] -> Bool
isSorted (x:y:zs) = (x <= y) && isSorted (y:zs)
isSorted _        = True
\end{minted}

However, since we use bind to pass permutations from \hinl{perm} 
to \hinl{isSorted}, each permutation is fully evaluated before 
it is determined whether the permutation is sorted. This leads to the complete 
evaluation of every permutation, which results in an inefficient program.

Similarly, when we consider the Curry example \hinl{head (1 : 
head [] : [])}, the strictness of our \hinl{MonadPlus} approach 
shows again. The corresponding Haskell expression is as follows.

\mint{Haskell}{hd [] >>= \x -> hd (1 : x : [])}

Here \hinl{hd :: MonadPlus a => [a] -> m a} is the lifted 
\hinl{head} function. Evaluating the expression in Haskell 
yields \hinl{mzero}, that is, no result, while Curry returns 
\hinl{1}. The reason is the definition of the bind operator. For 
example, the monad instance for lists defines bind as \hinl{xs 
>>= f = concatMap f xs}. In the expression above, this means that the pattern 
matching within \hinl{concatMap} evaluates 
\hinl{hd []} to \hinl{mzero} and thus returns 
\hinl{mzero}.

The strictness observed in both examples is the motivation for an alternative 
approach. The problem with the above implementations is that non-deterministic 
arguments of constructors need to be evaluated completely before the 
computation can continue. Therefore, we would like to be able to use 
unevaluated, non-deterministic computations as arguments of constructors.

As mentioned before, we can implement this idea by adapting all data types so 
that they may contain non-deterministic components.

\begin{minted}{Haskell}
data List m a = Nil | Cons (m a) (m (List m a))
\end{minted}

The list data type now has an additional argument \hinl{m} of 
type \hinl{* -> *} that represents a non-determinism monad. 
Instead of fixed constructors like \hinl{Choice}, the monad 
\hinl{m} determines the structure and evaluation strategy 
of the non-determinism effect. Two smart constructors 
\hinl{cons} and \hinl{nil} make handling the new 
list type more convenient.

\begin{minted}{Haskell}
nil :: Monad m => m (List m a)
nil = return Nil

cons :: Monad m => m a -> m (List m a) -> m (List m a)
cons x y = return (Cons x y)
\end{minted}

Adapting the permutation sort functions to the lifted data type requires us to 
replace \hinl{[]} with \hinl{List m}. However, 
this is not sufficient because the list itself can be the result of a 
non-deterministic computation. Therefore, an additional \hinl{m} 
is wrapped around every occurrence of \hinl{List}.

\begin{minted}{Haskell}
insert' :: MonadPlus m => m a -> m (List m a) -> m (List m a)
insert' mx mxs = cons mx mxs
  `mplus` mxs >>= \xs -> case xs of
                           Nil         -> mzero
                           Cons my mys -> cons my (insert' mx mys)

perm' :: MonadPlus m => m (List m a) -> m (List m a)
perm' ml = ml >>= \l ->
  case l of
    Nil -> nil
    Cons mx mxs -> insert' mx (perm' mxs)
\end{minted}

Whenever pattern matching occurred in the original definition, we now use bind 
to extract a \hinl{List} value. Since this only evaluates flat 
non-determinism and not non-determinism that occurs in the components, 
non-strictness is upheld as much as possible. 

All functions now take arguments of the same type they return. Thus, the 
definition of \hinl{sort} does not need bind in order to pass 
permutations to \hinl{isSorted}.

\begin{minted}{Haskell}
sort' :: MonadPlus m => m (List m Int) -> m (List m Int)
sort' xs = let ys = perm' xs in
  isSorted' ys >>= \sorted -> guard sorted >> ys
\end{minted}

We are now able to take advantage of \hinl{isSorted}'s non-strict
definition. The implementation generates permutations only if there is a chance 
that the permutation is sorted, that is, only recursive calls of 
\hinl{perm} that are demanded by \hinl{isSorted} 
are executed.

Data types with non-deterministic components solve the problem of 
non-strictness because each component can be evaluated individually, instead of 
forcing the evaluation of the whole term. Unfortunately, this leads to a 
problem. When unevaluated components are shared via Haskell's built-in sharing, 
computations, rather than results, are being shared. This means that the 
results can be different each time the computation is evaluated, which 
contradicts the intuition of sharing.

The solution to this problem is an explicit sharing combinator 
\hinl{share :: m a -> m (m a)} that allows sharing the results 
of a computation in a non-strict way. Here, \hinl{m} is a 
\hinl{MonadPlus} instance, similar to the monad used in the 
definition of the data type, that supports sharing. Thus, 
\hinl{share} takes a  computation and then returns a computation 
that returns the result, that is, the shared value. The reason for this nesting 
of monad layers is that, in short, the \hinl{share} combinator 
performs some actions that can be immediately executed by bind (the outer monad 
layer), while the inner monad layer should only be evaluated when needed. This 
is explained in more detail later.
With the explicit sharing operator we can adapt \hinl{perm'} to
share the generated permutations in order to achieve non-strictness in 
combination with sharing.

\begin{minted}{Haskell}
sort' :: MonadPlus m => m (List m Int) -> m (List m Int)
sort' xs = do ys <- share (perm' xs)
              sorted <- isSorted'
              guard sorted
              ys
\end{minted}

The \hinl{share} operator must satisfy certain laws, which we 
discuss in \autoref{sec:lawsOfSharing}. The implementation of 
\hinl{share} is 
subject of the next chapter.

\chapter{Call-Time Choice modelled in Haskell}
Based on the ideas presented in the last chapter, we now want to model 
call-time choice, that is, non-strictness, sharing and non-determinism, in 
Haskell. We still use \hinl{MonadPlus} to parameterize our 
programs. However, instead of, for example, using the list instance to make 
non-determinism visible, we define an effect functor that can express many 
different effects, including non-determinism and sharing. This approach, as 
introduced by \citep{wu2014effect}, will also be the base of the Coq 
implementation shown in \autoref{ch:callTimeChoiceCoq}. \todo{Definition effect}

For the implementation of call-time choice, we want to be able to express 
different effects within our programs. However, not every program contains 
effects. There are also \textit{pure} programs that have no side-effects 
besides the computation of a value. A data type that represents such programs 
could look as follows.

\begin{minted}{Haskell}
data Void a = Return a
\end{minted}

Here, \hinl{Void} means the absence of effects. 

If we consider programs that contain effects, also called \textit{impure}, 
like, non-determinism, a data type that represents such values could look like 
the following.

\begin{minted}{Haskell}
data ND a = Return a
          | Fail
          | Choice (ND a) (ND a)
\end{minted}

This data type also has a constructor to model pure values but in addition, 
there are constructors that represent failed computations and the 
non-deterministic choice between two values. We could go on and list data types 
that model many more effects but the question is: Is it possible to create a 
data type that, if appropriately instantiated, behaves like the original effect 
functor? This would allow us to represent programs with many different effects 
using one compact data type.

Answering this question requires abstracting the concrete form of effect 
functors into a general program data type. As we saw in the examples above, we 
need a way to represent pure values in a program. Therefore, the first 
constructor of our new program data type should be \hinl{Return 
a} for the type \hinl{a}, that is, the result type of the 
program. To model effects like non-determinism, the program is parameterized 
over effect functors of type \hinl{* -> *} that represent, for 
example, \hinl{Fail} and \hinl{Choice}. We call 
this argument \hinl{sig} because the signature of a program 
tells us which effects can occur. So far, programs are defined as 
\hinl{data Prog sig a = Return a}. In order do make use of the 
\hinl{sig} component, we need to add a constructor for impure 
operations. The \hinl{ND} data type shows us that effect 
functors can be defined recursively. Thus, the constructor for impure programs 
should be recursive, too, to be able to represent this structure.

\begin{minted}{Haskell}
data Prog sig a = Return a | Op (sig (Prog sig a))
\end{minted}

With this definition of \hinl{Prog}, we are able to represent 
the original functors by instantiating \hinl{sig} appropriately. 
For \hinl{Void}, we already have the 
\hinl{Return} constructor. Therefore, the data type we can use 
with \hinl{Prog} does not need a constructor anymore, that is, 
\hinl{data Void' a}. 

\begin{minted}{Haskell}
VoidProg a = Return a
           | Op (Void' (VoidProg a)) -- Void' has no constructors!
\end{minted}

The type \hinl{Prog Void'} now resembles the original type 
\hinl{Void} since the the \hinl{Op} constructor 
would require a value of type \hinl{Void'}, which we cannot 
construct.\footnote{It is possible to use \hinl{undefined} to 
create an impure value of type \hinl{Prog Void' a}. Since this 
is not possible in Coq, we do not consider this in the Haskell implementation.}
Only \hinl{Return} can be used to define values, similar to the 
original data type.

Similar to \hinl{Void'}, we can define a data type 
\hinl{Choice'} that represents \hinl{Choice} in 
combination with \hinl{Prog}. 

\begin{minted}{Haskell}
data Choice' p = Fail' | Choice' p p
\end{minted}

Again, we can omit the \hinl{Return} constructor because it is 
already part of the \hinl{Prog} data type. For the same reason, 
the type variable \hinl{a} has been replaced with the variable 
\hinl{p}, since \hinl{Choice} does not have 
values as arguments but rather programs that return values.

\begin{minted}{Haskell}
data ChoiceProg a = Return a
                  | Op (Choice' (ChoiceProg a))
\end{minted}

Since \hinl{Op} applies \hinl{sig} recursively, 
this yields the following type, which is equivalent to the original data type 
\hinl{Choice}.

\begin{minted}{Haskell}
data ChoiceProg a = Return a
                  | OpFail
                  | OpChoice (ChoiceProg a) (ChoiceProg a)
\end{minted}

We have found a way to model effect functors as instances of the data type 
\hinl{Prog}, which essentially models a tree with leafs, 
represented by the \hinl{Return} constructor, and branches that 
have the form defined by \hinl{sig}. \todo{Tree structure 
visualization?}

\section{Free Monads}
\begin{itemize}
\item What are free monads?
\item Why do we use free monads?
\end{itemize}
The data type \hinl{Prog} is better known as the \textit{free 
monad}. We saw in the previous chapter that \hinl{Free} can be 
used to model other data types. In addition, \hinl{Free} is a 
monad that can turn any functor into a monad.

We consider, for example, the type \hinl{Free One} where \hinl{data One a = 
One}. Here \hinl{a} is a phantom type that we need because \hinl{Free} expects 
a functor. The monad instance for \hinl{Free} is as follows.

\begin{minted}{Haskell}
data Free f a = Pure a | Impure (f (Free f a))

instance (Functor f, Applicative (Free f)) => Monad (Free f) where
  return = Pure
  Pure x >>= g = g x
  Impure fx >>= g = Impure (fmap (>>= g) fx)
\end{minted}

Since \hinl{One} has only a single, non-recursive constructor \hinl{One}, the 
only possible impure value is \hinl{Impure One}, whereas the usual 
\hinl{Return} constructor remains. If bind encounters the value \hinl{One}, the 
function \hinl{g} is distributed deeper into the term structure using 
\hinl{fmap}. Since \hinl{fmap One = One}, it becomes apparent that the monad 
constructed by \hinl{Free One} is the \hinl{Maybe} monad.
\todo{Visualization of fmap tree}

Since we want to model different effects in our program, the free monad makes 
writing programs easier by allowing monadic definitions without defining a 
separate monad instance for each effect.
\todo{Find more reasons for using free monads}
\section{Modelling Effects}
\begin{itemize}
\item Explanation of the Prog/sig infrastructure
\item ND and state effect implementation 
\end{itemize}

\section{Sharing}
\begin{itemize}
\item How can we implement simple sharing as an effect?
\item What about deep/nested sharing?
\item Examples (exRecList, ...)
\end{itemize}


\chapter{Call-Time Choice modelled in Coq}
\label{ch:callTimeChoiceCoq}
The goal of this chapter is to transfer the Haskell implementation of call-time 
choice to Coq. We begin with the data structure \hinl{Prog}, 
that is, the free monad, which allowed us to model programs with effects of 
type \hinl{sig} and results of type \hinl{a}.

\begin{minted}{Haskell}
data Prog sig a = Return a | Op (sig (Prog sig a))
\end{minted}

The definition in Coq looks very similar to the Haskell version, aside from 
renaming and the explicit constructor types.

\begin{minted}{Coq}
Inductive Free F A :=
| pure : A -> Free F A
| impure : F (Free F A) -> Free F A.
\end{minted}

However, the definition is rejected by Coq upon loading the file with the 
following error message.

\begin{minted}{Text}
Non-strictly positive occurrence of "Free" in "F (Free F A) -> Free F A".
\end{minted}

The reason for this error is explained in the next section.

\section{Non-strictly Positive Occurrence}
\begin{itemize}
\item What does non-strictly positive occurrence mean?
\item Motivation for usage of containers
\end{itemize}

In \autoref{sec:coqIntro}, we learned that Coq distinguishes between 
non-recursive definitions and functions that use recursion. The reason for this 
is that Coq checks functions for termination, which is an important part of 
Coq's proof logic. To understand why functions must always terminate in Coq, we 
consider the following function.

% https://www.di.ens.fr/~zappa/teaching/coq/ecole11/summer/lectures/lec9.pdf
\begin{minted}{Coq}
Fail Fixpoint loop (x : unit) : A := loop x.
\end{minted}

The function receives an argument \mintinline{Coq}{x} and calls itself with the 
same argument. Since this function obviously never terminates, the result type 
\mintinline{Coq}{A} is arbitrary. In particular, we could instantiate 
\mintinline{Coq}{A} with \mintinline{Coq}{False}, the false proposition. The 
value \mintinline{Coq}{loop tt : False} could be used to prove anything,
according to the principle of explosion. For this reason, Coq requires all 
recursive functions to terminate provably.

Returning to the original data type, what is link between 
\mintinline{Coq}{Free} and termination? It is well known that recursion can be 
implemented in languages without explicit recursion syntax by means of 
constructs like the Y combinator or the data type \mintinline{Coq}{Mu} for 
type-level recursion.

\begin{minted}{Coq}
Fail Inductive Mu A := mu : (Mu A -> A) -> Mu A.
\end{minted}

\mintinline{Coq}{Mu} is not accepted by Coq for the same reason as 
\mintinline{Coq}{Free}: non-strictly positive occurrence of the respective data 
type. The problematic property of non-strictly positive data types is that the 
type occurs on the left-hand side of a constructor argument's function type. 
This would allow general recursion and thus, as described above, make Coq's 
logic inconsistent.

In case of \mintinline{Coq}{Free}, the non-strictly positive occurrence is not 
as apparent as before because the constructors do not have functional 
arguments. However, \mintinline{Coq}{F} is being applied to 
\mintinline{Coq}{Free F A}. If \mintinline{Coq}{F} has a functional argument 
with appropriate types, the resulting type becomes non-strictly positive, as 
shown below.

\begin{minted}{Coq}
Definition Cont R A := (A -> R) -> R.

(* Free (Cont R) *)
Fail Inductive ContF R A :=
| pureC   : A -> ContF R A
| impureC : ((ContF R A -> R) -> R) -> ContF R A.
\end{minted}

In the type of \mintinline{Coq}{impureC} contains a non-strictly positive 
occurrence of \mintinline{Coq}{ContF R A}. Consequently, Coq rejects 
\mintinline{Coq}{Free} because it is not guaranteed that no instance violates 
the strict positivity requirement. Representing the \mintinline{Coq}{Free} data 
type therefore requires a way to restrict the definition to strictly positive 
data types. One approach to achieve this goal is described in the next section.

\section{Containers}

\begin{itemize}
\item How do containers work?
\item How do we translate effect functors into containers?
\end{itemize}

Containers are an abstraction of data types that store values, with the 
property that only strictly positive data types can be modelled as a container. 
This will allow us to define a version of \mintinline{Coq}{Free} that works 
with containers of type constructors instead of the type constructors itself. 
First, however, we will have a more detailed look at containers.

The first component of a container is the type \mintinline{Coq}{Shape}. A 
shape determines how the data type is structured, regardless of the stored 
values. For example, the  shape of a list is the same as the shape of Peano 
numbers: a number that  represents the length of the list, or rather the number 
of \hinl{Cons}/\hinl{Succ} applications. A pair, 
on the other hand, has only a single shape.

The second component of a container is a function \mintinline{Coq}{Pos : Shape 
-> Type} that gives each shape a type that represents the positions within the 
shape. In the example of pairs, the shape has two positions, the first and 
second component. Each element of a list is a position within the shape. 
Therefore, the position type for lists with length $n$ is natural numbers 
smaller than $n$. Peano numbers do not have elements and therefore, the 
position type for each shape is empty.

Containers can be extended by a function that maps all valid positions to 
values. Since the position type depends on a concrete shape, the definition in 
Coq is quantified universally over values of type \mintinline{Coq}{Shape}.

\begin{minted}{Coq}
Inductive Ext Shape (Pos : Shape -> Type) A := 
  ext : forall s, (Pos s -> A) -> Ext Shape Pos A.
\end{minted}

The extension of a container models the concrete data type.

\section{Modelling Effects}
\begin{itemize}
\item In which ways is the Coq implementation simplified, compared to Haskell?
\item How does the adapted Prog/sig infrastructure work?
\item How do we translate recursive functions?
\end{itemize}

\section{Sharing}
\label{sec:lawsOfSharing}
\begin{itemize}
\item Laws of sharing
\end{itemize}

\chapter{Curry Programs modelled in Coq}
\begin{itemize}
\item Can we use the Coq model of call-time choice to prove properties about 
actual Curry programs?
\end{itemize}

\chapter{Conclusion}


% Literatur
\bibliographystyle{plainnat}
\bibliography{Mathesis} % Datei: seminar.bib
% Anhang
\appendix
\end{document}